\documentclass[14pt,a4paper]{extarticle}
\usepackage[T2A]{fontenc}
\usepackage[utf8]{inputenc}
\usepackage[english,russian]{babel}
\usepackage{lmodern}
\usepackage{geometry}
\usepackage{array}
\usepackage{tabularx}
\usepackage{longtable}
\usepackage{ltablex}
\usepackage{booktabs}
\usepackage{hyperref}
\usepackage{indentfirst}
\usepackage{microtype}
\usepackage{graphicx}
\usepackage{setspace}
\usepackage{listings}
\usepackage{graphicx}
\usepackage{caption}
\usepackage{float}
\usepackage{tocloft}
\usepackage{titlesec}
\usepackage{newtxtext,newtxmath}

\titleformat{\section}        % форматируем \section
  {\normalfont\bfseries}      % стиль: обычный + жирный
  {\thesection.}              % номер секции с точкой
  {2mm}                       % расстояние между номером и текстом
  {}                           % перед заголовком можно добавить что-то
\titlespacing*{\section}{1.5cm}{*3}{*1}  % {отступ слева}{отступ сверху}{отступ снизу}

\titleformat{\subsection}        % форматируем \section
  {\normalfont\bfseries}      % стиль: обычный + жирный
  {\thesubsection.}              % номер секции с точкой
  {2mm}                       % расстояние между номером и текстом
  {}                           % перед заголовком можно добавить что-то
\titlespacing*{\subsection}{1.5cm}{*3}{*1}  % {отступ слева}{отступ сверху}{отступ снизу}

\captionsetup{justification=centering}
\renewcommand{\ttdefault}{cmtt}
\keepXColumns
\setlength{\parindent}{3em}
\newgeometry{left=2.5cm,right=2.5cm,top=3cm,bottom=2.5cm}

\renewcommand{\contentsname}{\hspace*{4em}{\smallСодержание}}
\setlength{\cftbeforesecskip}{0.5em}
\renewcommand{\cftsecleader}{\cftdotfill{\cftdotsep}}
\renewcommand{\cftsecaftersnum}{.\hspace{2mm}}

\title{Отчёт по "Управление ИТ-проектами"}
\author{Выполнил студент группы М-РИТ-25-1, Станиславчук Сергей}

\begin{document}

\setstretch{1.5}
\noindent
\begin{minipage}[t]{0.01\textwidth}
    \includegraphics{~/programmer/++Programmer/5_1/LSTU.jpg}
\end{minipage}%
\begin{minipage}[c]{1.1\linewidth}
    \centering
    {\small
    МИНИСТЕРСТВО НАУКИ И ВЫСШЕГО ОБРАЗОВАНИЯ РФ \\[1em]
    ФЕДЕРАЛЬНОЕ ГОСУДАРСТВЕННОЕ БЮДЖЕТНОЕ \\
    ОБРАЗОВАТЕЛЬНОЕ УЧРЕЖДЕНИЕ ВЫСШЕГО ОБРАЗОВАНИЯ \\
    «ЛИПЕЦКИЙ ГОСУДАРСТВЕННЫЙ ТЕХНИЧЕСКИЙ УНИВЕРСИТЕТ»
    }
\end{minipage}%

\vspace{3em}

\noindent
\hspace{1em}%
\begin{minipage}[t]{0.8\textwidth}
\textbf{Институт} \hspace{2em} компьютерных наук \\
\textbf{Кафедра} \hspace{2em} автоматизированных систем управления
\end{minipage}


\vspace{6em}

\begin{center}
\textbf{Практическая работа №4} \\
по дисциплине \\ <<Управление ИТ-проектами>>
\end{center}

\vspace{6em}

\noindent
\begin{tabular*}{\textwidth}{@{\extracolsep{\fill}} l c r }
    Студент М-РИТ-25-1 & \makebox[4cm]{\hrulefill} & Станиславчук С.М. \\[2em]
Руководитель & \\
К.т.н., доцент & \makebox[4cm]{\hrulefill} & Ведищев В.В. \\
\end{tabular*}

\vfill
\begin{center}
    Липецк, 2026
\end{center}

\thispagestyle{empty}

\newpage
\tableofcontents
\newpage

\section{Разработка общей концепции плана проекта}

Общая концепция плана проекта определяет структуру работ по созданию системы корпоративного обмена сообщениями, включая последовательность этапов, содержание работ и используемые ресурсы.

\subsection{Этапы выполнения проекта}

Процесс разработки системы включает следующие этапы:

\begin{itemize}
\item \textbf{Сбор и анализ требований} — изучение процессов корпоративной коммуникации, определение потребностей пользователей и ограничений системы;
\item \textbf{Формализация и спецификация требований} — оформление требований в виде структурированного документа (спецификации);
\item \textbf{Проектирование системы} — разработка архитектуры, определение компонентов (клиент, сервер, протоколы взаимодействия);
\item \textbf{Реализация системы} — разработка программных модулей, реализующих функции обмена сообщениями, голосовой связи и передачи данных;
\item \textbf{Тестирование} — проверка корректности работы системы, выявление и устранение ошибок, верификация требований;
\item \textbf{Внедрение и эксплуатация} — развёртывание системы в корпоративной сети и её использование в реальных условиях.
\end{itemize}

\begin{table}[h!]
\centering
\begin{tabular}{|p{5cm}|p{5cm}|p{4cm}|} \\
\hline
\textbf{Этап} & \textbf{Содержание работ} & \textbf{Ресурсы} \\
\hline
Сбор требований & Анализ процессов, выявление потребностей пользователей & Аналитик, эксперт по ИБ \\
\hline
Спецификация & Формализация требований, подготовка документации & Аналитик, разработчик \\
\hline
Проектирование & Разработка архитектуры системы & Бэкенд-разработчик, сетевой разработчик \\
\hline
Реализация & Программирование клиентской и серверной части & Фронтенд, бэкенд разработчики \\
\hline
Тестирование & Проверка функциональности и безопасности & Тестировщик, разработчик \\
\hline
Эксплуатация & Развёртывание и поддержка системы & Системный администратор \\
\hline
\end{tabular}
\caption{План работ и ресурсы проекта}
\end{table}

\subsection{Детальные требования}

\textbf{Сетевые протоколы передачи данных}

Для реализации обмена данными в системе используются протоколы, обеспечивающие работу в реальном времени:
\begin{itemize}
\item \textbf{WebSocket} — для двустороннего обмена текстовыми сообщениями и служебными данными (события, статусы пользователей);
\item \textbf{UDP} — для передачи голосового трафика с минимальными задержками;
\item \textbf{RTP (Real-time Transport Protocol)} — для структурирования аудиопотока;
\item \textbf{HTTP/HTTPS} — для передачи файлов и выполнения вспомогательных операций (авторизация, загрузка ресурсов).
\end{itemize}

\textbf{Шифрование и защита данных}

Система должна обеспечивать защиту данных на всех этапах передачи:
\begin{itemize}
\item использование \textbf{TLS} для защиты WebSocket и HTTP-соединений;
\item применение \textbf{DTLS} для защиты UDP-трафика;
\item безопасное хранение паролей (хэширование \textbf{bcrypt});
\item реализация механизмов аутентификации и авторизации (токены, сессии);
\end{itemize}

\textbf{Требования к задержкам и производительности}

\begin{itemize}
\item задержка доставки текстовых сообщений — не более 100–200 мс в пределах локальной сети;
\item задержка голосовой связи — не более 150 мс (end-to-end);
\item поддержка одновременной работы большого количества пользователей (десятки–сотни в рамках сети);
\item эффективное использование пропускной способности сети;
\item минимизация потерь пакетов и механизм повторной передачи при необходимости.
\end{itemize}

\textbf{Требования к пользовательскому интерфейсу}

\begin{itemize}
\item интерфейс должен быть интуитивно понятным и минималистичным;
\item наличие чатов (личные и групповые);
\item отображение списка пользователей и их статусов (онлайн/оффлайн);
\item поддержка голосовых каналов;
\item уведомления о новых сообщениях и событиях;
\item возможность отправки файлов через drag-and-drop;
\item высокая отзывчивость интерфейса.
\end{itemize}

\textbf{Масштабируемость системы}

\begin{itemize}
\item модульная архитектура (разделение на сервисы: аутентификация, сообщения, голос);
\item возможность горизонтального масштабирования серверной части;
\item оптимизация хранения и передачи данных (кэширование, сжатие).
\end{itemize}

Представленные требования ориентированы на построение системы, аналогичной по принципам работы современным платформам обмена сообщениями, но адаптированной для использования в корпоративной среде с повышенными требованиями к безопасности и производительности.

\newpage
\section{Структурная декомпозиция работ}

Структурная декомпозиция проекта представляет собой разбиение процесса разработки системы на иерархию работ, упорядоченных от конечного результата к элементарным задачам. Такой подход позволяет определить состав работ, их взаимосвязи и уровень детализации, необходимый для планирования и контроля выполнения проекта.

В корне структуры находится конечный результат — функционирующая система корпоративного обмена сообщениями, обеспечивающая текстовую и голосовую связь в реальном времени. Достижение данного результата обеспечивается выполнением совокупности взаимосвязанных этапов, каждый из которых, в свою очередь, декомпозируется на более простые задачи.

\subsection*{Иерархия работ}

\begin{itemize}
\item \textbf{1. Готовая система}
\begin{itemize}
\item \textbf{1.1 Сбор и анализ требований}
\begin{itemize}
\item 1.1.1 Анализ процессов коммуникации
\item 1.1.2 Определение пользователей и ролей
\item 1.1.3 Выявление ограничений системы
\end{itemize}

    \item \textbf{1.2 Спецификация требований}
    \begin{itemize}
        \item 1.2.1 Формализация требований
        \item 1.2.2 Согласование с заинтересованными сторонами
        \item 1.2.3 Утверждение документа
    \end{itemize}

    \item \textbf{1.3 Проектирование системы}
    \begin{itemize}
        \item 1.3.1 Разработка архитектуры
        \item 1.3.2 Выбор протоколов и технологий
        \item 1.3.3 Проектирование базы данных
        \item 1.3.4 Проектирование интерфейса
    \end{itemize}

    \item \textbf{1.4 Реализация системы}
    \begin{itemize}
        \item 1.4.1 Разработка серверной части
        \begin{itemize}
            \item 1.4.1.1 Реализация API
            \item 1.4.1.2 Реализация обработки сообщений
            \item 1.4.1.3 Реализация голосового модуля
        \end{itemize}
        \item 1.4.2 Разработка клиентской части
        \begin{itemize}
            \item 1.4.2.1 Интерфейс пользователя
            \item 1.4.2.2 Работа с сетью
            \item 1.4.2.3 Обработка событий
        \end{itemize}
    \end{itemize}

    \item \textbf{1.5 Тестирование}
    \begin{itemize}
        \item 1.5.1 Функциональное тестирование
        \item 1.5.2 Тестирование производительности
        \item 1.5.3 Тестирование безопасности
    \end{itemize}

    \item \textbf{1.6 Внедрение и эксплуатация}
    \begin{itemize}
        \item 1.6.1 Развёртывание системы
        \item 1.6.2 Настройка окружения
        \item 1.6.3 Поддержка пользователей
    \end{itemize}
\end{itemize}

\end{itemize}

Данная декомпозиция отражает последовательный переход от общей цели проекта к конкретным элементарным задачам, выполнение которых обеспечивает достижение конечного результата.
\end{document}
